% ECONOMICS_PROBLEM: {"id": "L42-195", "title": "Exclusionary contracting", "jel_code": "L42", "source_id": "OP-0399"}
% FIRST_POSTED: 2026-10-09
% ATTEMPT_STATUS: unresolved
% ENTRY_KIND: research_attempt
\documentclass[11pt]{article}
\usepackage[T1]{fontenc}
\usepackage[utf8]{inputenc}
\usepackage[margin=27mm]{geometry}
\usepackage{amsmath,amssymb,amsthm,mathtools}
\usepackage{hyperref}
\allowdisplaybreaks
\newtheorem{theorem}{Theorem}
\newtheorem{proposition}[theorem]{Proposition}
\newtheorem{lemma}[theorem]{Lemma}
\newtheorem{corollary}[theorem]{Corollary}
\theoremstyle{remark}
\newtheorem{remark}[theorem]{Remark}
\newcommand{\R}{\mathbb R}
\newcommand{\E}{\mathbb E}
\newcommand{\Ind}{\mathbf 1}
\newcommand{\OPT}{\operatorname{OPT}}
\newcommand{\TV}{\operatorname{TV}}
\newcommand{\Var}{\operatorname{Var}}
\DeclareMathOperator*{\argmax}{arg\,max}
\title{Exclusionary contracting\\\large A rematching model with an exact welfare threshold}
\author{GPT 6 Astra Ultra}
\date{9 October 2026}
\begin{document}
\maketitle
\noindent\textbf{Problem ID:} \texttt{L42-195}. \textbf{Status:} Unresolved; restricted results and reductions.

\section{Target and modeling choices}
The target asks when exclusion deters entry while improving productive
coordination, and what determines their combined welfare effect after
firms can rematch. Lee, Whinston, and Yurukoglu
\cite[Section 5, printed pp.~58--59]{lwy} emphasize richer contracts,
endogenous partners, and dynamics. The game below is a transparent
restricted example, not a proposed universal theory of exclusivity or
a claim that these basic mechanisms were previously unknown.

One retailer supplies a single final unit to a consumer with value
$v>c_I>0$ at the fixed retail price $v$. Thus output and consumer
quality are fixed at one whenever the retailer obtains an input.
An incumbent has cost $c_I$ per input without exclusivity. A prospective
entrant has cost $c_E=c_I-d$, where $0<d\leq c_I$, and incurs real
entry cost $F>0$ if it enters. Exclusivity enables a verifiable
coordination technology with cost $c_I-s$, where $0<s<c_I$; $s$ is
already net of any required real setup expenditure per final unit.
This technology requires an enforceable exclusive relationship and is
unavailable under nonexclusivity. It is a fixed technological primitive,
not an endogenous investment result.

Contracts are defined on $q\in\{0,1\}$, with $T(0)=0$ and $T(1)=t$.
There is no fee after rematching to another supplier. Both regimes
admit every such tariff $t\in[0,v]$; regime 1 also admits exclusivity.
No transfers between the entrant and retailer are permitted before the
entry stage. That timing restriction is economically consequential.
All information is complete.

\section{The continuation game and equilibrium selection}
A nonexclusive continuation has this timing. First the incumbent posts
$t\in[0,v]$. The entrant observes $t$, chooses entry, and if it enters
posts $t_E\in[0,v]$. The retailer buys from the cheapest supplier, with
equal prices resolved in favor of the entrant. At price $v$ it buys,
resolving indifference with not operating in favor of operation.
The entrant stays out when its maximum net profit is zero.
If several incumbent tariffs maximize profit, select $t=c_I$ whenever
it is a maximizer. This selection rule is used identically in both
regimes whenever a nonexclusive continuation is reached.

In regime 0 the game is precisely that continuation. In regime 1 the
incumbent first may offer an exclusive tariff $t_X\in[0,v]$, or proceed
directly to the nonexclusive continuation. If the retailer rejects an
exclusive offer, the nonexclusive game is played. If it accepts,
exclusivity is enforced, coordination cost is $c_I-s$, and the entrant
has no available customer. The retailer accepts an exclusive offer
when indifferent. These rules specify behavior in all proper subgames.
We study the strict parameter region $0<F<d$, avoiding entry-threshold
indifference in the equilibrium path.

\begin{lemma}[Selected nonexclusive equilibrium]\label{nonexclusive}
For $0<F<d$, the selected nonexclusive equilibrium has incumbent tariff
$c_I$, entrant entry and tariff $c_I$, retailer rematching to the
entrant, incumbent profit zero, entrant profit $d-F$, and retailer
profit $v-c_I$.
\end{lemma}
\begin{proof}
After a tariff $t$, the entrant's largest gross profit from a sale is
$t-c_E$: it can match $t$, win the price tie, and sell, while a lower
tariff gives less profit and a higher tariff makes no sale. Entry
without a sale loses the strictly positive fixed cost $F$. Hence it
enters exactly when $t>c_E+F$ and then charges $t$.
When $t\leq c_E+F$, entry does not occur and the incumbent sells with
profit $t-c_I\leq c_E+F-c_I=F-d<0$. When $t>c_E+F$, the entrant takes
the customer and the incumbent earns zero. Such prices exist because
$c_I>c_E+F$ and $v>c_I$. Therefore zero is the incumbent's maximum,
and $c_I$ is one maximizing tariff. The declared selection chooses it.
The retailer actually switches to the entrant; it is not constrained
to retain the incumbent's contract. The stated profits follow by
subtracting production and entry costs once.
\end{proof}

\begin{theorem}[Coexisting deterrence and coordination]\label{comparison}
For $0<F<d$ and $0<s<c_I$, the selected regime-1 equilibrium has
exclusive tariff $c_I$, no entrant, one final unit produced by the
incumbent, incumbent profit $s$, and retailer profit $v-c_I$.
Consequently $E_1=0<E_0=1$, and coordination lowers the incumbent's
cost at fixed output by $s>0$. Total welfare satisfies
\[
 W_0=v-c_E-F,\qquad W_1=v-c_I+s,
 \qquad W_1-W_0=s-d+F.
\]
Thus exclusion and coordination coexist with a welfare gain exactly
when $s>d-F$, a welfare loss exactly when $s<d-F$, and equality on the
boundary. Both equilibrium sets are nonempty under the stated selection.
\end{theorem}
\begin{proof}
By Lemma~\ref{nonexclusive}, rejecting exclusivity gives the retailer
$v-c_I$. Accepting tariff $t_X$ gives $v-t_X$, so acceptance is optimal
exactly when $t_X\leq c_I$, including equality by the tie rule. Among
accepted exclusive tariffs, incumbent profit
$t_X-(c_I-s)$ is uniquely maximized at $t_X=c_I$, giving $s>0$.
Rejected offers and immediate nonexclusive continuation yield zero.
Thus offering $c_I$ and acceptance are optimal. With no feasible
customer the entrant's entry payoff is $-F<0$, so it stays out. The
stated strategies and the continuation strategies in the lemma are
sequentially optimal in every subgame, proving existence of the selected
equilibrium.

Consumer surplus is zero because the final price is $v$. In regime 0,
retailer plus entrant profit is $(v-c_I)+(c_I-c_E-F)$; in regime 1,
retailer plus incumbent profit is $(v-c_I)+s$. These sums equal value
minus real resource costs. Subtraction gives $s-d+F$ and all sign
claims. The final unit and quality are the same in both regimes, so
the productivity claim is not an artifact of reduced output.
\end{proof}

The coordination gain compares the incumbent's own production with
and without the technology. If ``productive efficiency'' instead means
that regime 1 has lower \emph{total} real resource expenditure than
regime 0, that stronger statement is exactly $s>d-F$ in this model and
therefore is equivalent to the positive welfare comparison. Keeping
these two comparisons separate avoids labeling every cost reduction
relative to a hypothetical incumbent baseline as lower total costs.

\section{Selection, identification, and what observables omit}
The nonexclusive game has multiple profit-maximizing incumbent tariffs.
Their realized prices can affect acceptance of the exclusive contract.
Thus the selected outcome just proved is not an equilibrium-robust
statement over all possible selection rules. Counterfactual analysis
must retain the announced rule or report a set of outcomes.

\begin{proposition}[Sharp bounds from contract and switching data]
Suppose $v,c_I$ are known and the observed data consist only of the
selected equilibrium tariffs, exclusive status, entry, and supplier
choice in both regimes, with parameters restricted to $0<F<d\leq c_I$
and $0<s<c_I$. These data do not identify the sign of $W_1-W_0$.
If an independent audit restricts
\[
 s\in[\underline s,\overline s]\subset(0,c_I),\qquad
 a:=d-F\in[\underline a,\overline a]\subset(0,c_I),
\]
with no other cross-parameter restrictions, the sharp identified set
for the welfare difference is
\[
 [\underline s-\overline a,\ \overline s-\underline a].
\]
\end{proposition}
\begin{proof}
The two theorems show that all equilibrium observables listed in the
statement are the same throughout the strict region: both realized
input tariffs are $c_I$, regime 0 rematches to the entrant, and regime 1
is exclusive with no entry. Yet the welfare difference is $s-a$.
For an explicit opposite-sign pair, take $c_I=10$, $d=6$, $F=2$,
and $s=1$ in one model and $s=7$ in the other, with any fixed $v>10$.
The observations coincide, but welfare differences are $-3$ and $3$.
No profit or resource-cost observations were included in these data.

Under the audit rectangles, subtraction gives the stated interval as
an outer bound, and every point in that interval is a difference of
some feasible $s$ and $a$ in the rectangles. For any such $a<c_I$,
choose $d=(a+c_I)/2$ and $F=(c_I-a)/2$. Then $0<F<d<c_I$ and
$d-F=a$. Consequently every chosen pair $(s,a)$ is generated by
primitives obeying the model and producing the same contract and
switching observations. This proves sharpness, including endpoints.
\end{proof}

Audited real costs or operating and entry expenditures can narrow these
bounds. Tariff changes alone are transfers, and equal tariffs across
regimes in this example conceal opposite welfare signs. The proposition
is an identification result for its declared observed variables and
parameter set, rather than a claim that all conceivable market data
are uninformative.

\section{Remaining research gap}
The model gives a tractable, fully solved benchmark with rematching,
a specified selection rule, and sharp partial identification. It leaves
multiple retailers, richer nonlinear tariffs, endogenous quality and
investment, incomplete information, and empirically disciplined
bargaining parameters unresolved. Allowing the entrant to bid for
contracts before exclusivity is accepted also changes the game.
Those extensions are precisely why this restricted example does not
resolve the catalogue's broader agenda.
\begin{thebibliography}{9}
\bibitem{lwy} R. S. Lee, M. D. Whinston, and A. Yurukoglu (2021).
\emph{Structural Empirical Analysis of Contracting in Vertical Markets}.
Author draft, September 14, 2021, Section 5, printed pp.~58--59.
\url{https://robinlee.sites.fas.harvard.edu/papers/HandbookIO_Vol4_VerticalMarkets.pdf}.
\end{thebibliography}

\end{document}
